\section{Introduction}

Rapid advancement and scaling of large language models (LLMs) have led to the emergence of intelligent agent systems utilizing LLMs as cognitive cores to perceive environments, make autonomous decisions, and interact dynamically with external systems across various domains such as web automation~\citep{zhou2023webarena,drouin2024workarena}, software engineering~\citep{jimenez2023swe}, and embodied robotics~\cite{shridhar2020alfworld}. Despite demonstrating remarkable versatility, LLM-based agents inherit a critical limitation from their underlying neural architecture---hallucination, the phenomenon of generating outputs inconsistent with their contextual input~\citep{ji2023survey}. 
Unlike conventional natural language generation (NLG) tasks, hallucinations in agent settings translate directly into undesirable actions, amplifying real-world risks due to model autonomy and human oversight. 
% While agents may perform wrong actions due to general factors like domain knowledge gaps, flawed planning, or limited memory, hallucinated actions are a distinct subset driven by unfaithful interpretations of the contextual input.

Hallucination phenomena have been observed across various agent benchmarks. In TheAgentCompany~\citep{xu2024theagentcompany}, a notable instance is the agent unintentionally discloses sensitive user credentials due to hallucinated context switching. In another case from SWE-Bench~\cite{jimenez2023swe,yang2024swe}, the agent includes a non-existent package to forcibly resolve the bug. While such occasions are reported in agent studies, they remain scattered. Without a unified collection of test cases to gain deeper understanding of when and why hallucinations occur, it becomes increasingly difficult to develop reliable mitigation strategies or benchmark evaluations. In this work, our aim is to bridge this critical gap.

% \YY{CDiscuss with Figure jointly}.

Assessing hallucinations in interactive agent environments poses unique methodological challenges compared to traditional benchmarks, which typically focus on single-turn QA scenarios such as TruthfulQA~\cite{lin2021truthfulqa} or HaloGEN~\cite{ravichander2025halogen}. First, the output of an agent cannot be directly verified using external factuality scorers like FACTScore~\cite{min2023factscore}. Whether an action is hallucinatory depends heavily on the surrounding context, making verification more complex. Second, agents operate in diverse environments such as web-based interfaces~\cite{deng2023mind2web,yao2022webshop,zhou2023webarena,drouin2024workarena,lu2024weblinx,xu2024theagentcompany}, operating systems~\cite{xie2024osworld,wu2024copilot}, or Android applications~\cite{rawles2024androidworld, bai2024digirl}. These settings complicate the design of a unified evaluation framework~\cite{yehudai2025survey}. Finally, the stochastic nature of branching decisions makes agentic behavior less deterministic and harder to reproduce reliably.

To tackle these challenges, we introduce a unified benchmarking framework \textsc{MIRAGE-Bench} that \emph{categorizes}, \emph{elicits}, and \emph{evaluates} hallucinations in interactive LLM-agent scenarios:

$\cdot$ To \textit{categorize}, we extend from prior work that distinguishes hallucinations as “unfaithful to context” versus “unfaithful to factual knowledge”~\citep{ji2023survey,huang2025survey,tonmoy2024comprehensive}. We refine this taxonomy for agentic settings into three categories: a) \textbf{Unfaithful to task instructions} --- the agent executes an action that violates task goals, exceeds constraints, or relies on unstated user intents;
b) \textbf{Unfaithful to interaction history} --- the agent repeats completed steps, ignores prior outcomes, or otherwise contradicts its own action-observation trajectory, 
c) \textbf{Unfaithful to environment observations} --- the agent hallucinates elements or properties absent from the environment, such as clicking nonexistent buttons or assuming unreached states.

$\cdot$ To \textit{elicit} hallucinations, we perform an extensive audit of existing agent benchmarks~\citep{zhou2023webarena,xie2024osworld,drouin2024workarena,xu2024theagentcompany,jimenez2023swe,yao2024taubenchbenchmarktoolagentuserinteraction} to identify \emph{hallucination-prone risk settings}. These risk settings are further taxonomized and used to synthesize new test scenarios, enabling broader and more systematic coverage. Because interactive trajectories often branch stochastically, it is difficult to reliably reproduce hallucination behavior. To address this, we adopt a \textbf{contextual snapshot strategy} in which each test case freezes the agent's state immediately before a potential hallucination point. The LLM is then prompted to generate the next action within a controlled, deterministic context.

$\cdot$ To \textit{evaluate} hallucinations, we adopt a fine-grained-level \emph{LLM-as-a-Judge} paradigm~\citep{gu2024survey,zhuge2024agent} to design risk risk-setting-specific prompts that guide a separate judge model to semantically connect the motivation and intention behind a candidate action with the agent’s surrounding context, and evaluate its faithfulness accordingly. This method enables scalable assessment while preserves fidelity to nuances in agentic hallucination.

To the best of our knowledge, \textsc{MIRAGE-Bench} is the first framework to offer a systematic approach to studying hallucinations in interactive LLM agents, our paper contributes by:
\begin{enumerate}[leftmargin=*,topsep=0.5ex,itemsep=-0.5ex,partopsep=0.75ex,parsep=0.75ex,partopsep=0pt,wide,labelindent=0pt]
    \item Proposing a \textbf{unified taxonomy} of hallucinations in LLM-agent settings, distinguishing three core types: unfaithful to task instructions, unfaithful to execution history, and unfaithful to environment observations.
    \item Introducing a \textbf{snapshot-based elicitation strategy} that enables reliable reproduction of hallucination behavior by freezing interaction histories and prompting next-step predictions in a controlled environment.
    % \item 
    \item Releasing a \textbf{benchmark and evaluation toolkit} to support future research on diagnosing and mitigating hallucinations in interactive agent environments, featuring a scalable LLM-as-a-Judge framework with risk-aware prompts for high-fidelity, domain-general verification.
    \item Providing \textbf{empirical insights} that challenge common assumptions: hallucinations persist even in state-of-the-art LLMs based agents, the gap between open and closed models is narrower than expected, and hallucinations often follow recurring patterns (e.g., fabricating buttons, assuming prior success, or inventing details)---all posing distinct risks in agentic settings beyond NLG.
\end{enumerate}